\section{Introduction}
\label{introduction}

\subsection{Background and Motivation}
\label{intro:background}

Requirements determine what a software project is expected to deliver, but natural language requirements frequently include goals that are unrealistic under available resources, deployment environments, performance budgets, or organizational constraints.
Recent work has highlighted the central role of requirements in AI-driven software engineering~\cite{robinsonRequirementsAreAll2025} and the need to reduce ambiguity while preserving stakeholder readability in requirement specification~\cite{darifControlledNaturalLanguage2026}.
This paper focuses on \emph{requirement feasibility}: whether a stated requirement is realistic, implementable, and deployable under explicit or implicit constraints.

\subsection{Challenges of Requirement Feasibility Analysis}
\label{intro:challenges}

Feasibility analysis is difficult because many infeasible requirements are not syntactically invalid.
They may look precise while hiding resource limits, temporal assumptions, hardware dependencies, data availability assumptions, or deployment restrictions.
LLM-based requirement generation can amplify this issue by producing fluent but unrealistic requirements whose feasibility has not been checked against project knowledge.

\subsection{Multi-Agent Requirement Verification}
\label{intro:multi_agent}

We study a lightweight multi-agent verification workflow in which agents play complementary analysis roles: requirement interpretation, constraint verification, and resource/deployment feasibility checking.
The goal is not to build a large platform, but to decompose feasibility verification into focused reasoning steps that can be inspected and corrected by humans.

\subsection{Contributions}
\label{intro:contributions}

This paper makes the following planned contributions:
\begin{itemize}[leftmargin=*]
  \item A problem formulation for knowledge-guided requirement feasibility analysis and requirement realism verification.
  \item A multi-agent workflow for identifying unrealistic requirements, infeasible constraints, and missing feasibility assumptions.
  \item A constraint reasoning scheme covering resource, performance, deployment, and temporal feasibility.
  \item An experimental design for evaluating feasibility detection accuracy, constraint verification accuracy, unrealistic requirement detection, and human verification reduction.
\end{itemize}
